\section*{Capítulo \ref{ch:b2:field-lab-methods-2} --- Métodos de campo e de laboratório II}
% numbers from the Book 4 second-writer scratch calculator calc/ch31.py

\begin{solution}{exo:b2:field-lab-methods-2:1}
(a) (i) \omterm{def:b2:field-lab-methods-2:textures}{Textura ofítica}; (ii) \omterm{def:b2:field-lab-methods-2:textures}{textura granoblástica}; (iii) um \omterm{def:b2:field-lab-methods-2:porphyroblast}{poiquiloblasto}, um \omterm{def:b2:field-lab-methods-2:porphyroblast}{porfiroblasto} cheio de
inclusões; (iv) uma \omterm{def:b2:field-lab-methods-2:porphyroblast}{coroa}, uma orla de reação em volta de uma granada que estava se decompondo. (b) Por
\eqref{eq:b2:field-lab-methods-2:lami}, $\gamma_1/\gamma_2 = \sin\ang{104}/\sin\ang{128} =
0.9703/0.7880 = 1.23$: esse limite puxa \qty{23}{\percent} mais forte que os outros dois.
\end{solution}

\begin{solution}{exo:b2:field-lab-methods-2:2}
(a) $Q + A + P = 80$: Q $= 100 \times 25/80 = 31.3$, A $= 18.8$, P $= 50.0$. (b) P/(A+P) $=
40/55 = \qty{72.7}{\percent}$, entre 65 e 90 com Q entre 20 e 60: um granodiorito.
(c) $M' = 12 + 8 = \qty{20}{\percent}$, abaixo de 35: leucocrático.
\end{solution}

\begin{solution}{exo:b2:field-lab-methods-2:3}
(a) $\sigma = \sqrt{0.12 \times 0.88/800} = 0.0115$, isto é, $12.0 \pm 1.1$ pontos percentuais, um
erro relativo de $0.0115/0.12 = \qty{9.6}{\percent}$. (b) Quatro vezes mais pontos o reduzem à metade: $0.0057$.
\end{solution}

\begin{solution}{exo:b2:field-lab-methods-2:4}
(a) A perda é de \qty{0.0288}{g}, e $0.0288/0.8000 = \qty{3.60}{\percent}$.
(b) $0.7712/8.7712 = 0.0879$.
(c) Pó calcinado: $5.20/0.08792 = \qty{59.1}{\percent}$; pó original:
$59.14 \times (1 - 0.0360) = \qty{57.0}{\percent}$.
\end{solution}

\begin{solution}{exo:b2:field-lab-methods-2:5}
(a) $\sigma = 0.10 \times 0.02 = 0.002$ e $n = 0.02 \times 0.98/0.002^2 = 4900$ pontos.
(b) $\sigma = 0.03$ e $n = 0.30 \times 0.70/0.03^2 = 233.3$, ou seja, 234 pontos. (c) Para um erro relativo
fixo, $n = (1-p)/(p\,\varepsilon^2)$ cresce como $1/p$: o mineral acessório exige 21 vezes mais
pontos, mais do que uma lâmina costuma oferecer.
\end{solution}

\begin{solution}{exo:b2:field-lab-methods-2:6}
(a) RGM-1: $72.89/73.45 = \qty{99.2}{\percent}$; BCR-2: $54.17/54.1 = \qty{100.1}{\percent}$.
(b) RGM-1: $0.19/\sqrt{20} = 0.042$, e $-0.56/0.042 = -13.2$; BCR-2: $0.13/\sqrt{20} = 0.029$, e
$0.07/0.029 = 2.4$. (c) O RGM-1 tem viés: sua média se afasta do valor recomendado 13 vezes seu próprio
desvio-padrão, muito mais do que o acaso permite (se o valor recomendado estiver certo). O BCR-2 está no
limite do que o acaso dá. O RGM-1 não é menos preciso: seu desvio-padrão relativo,
$0.19/72.89 = \qty{0.26}{\percent}$, é próximo do do BCR-2, \qty{0.24}{\percent}. O problema é a
veracidade, não a precisão.
\end{solution}

\begin{solution}{exo:b2:field-lab-methods-2:7}
(a) $M = 39.098 + 3 \times 24.305 + 26.982 + 3 \times 28.085 + 12 \times 15.999 + 2 \times 1.008 =
\qty{417.25}{g.mol^{-1}}$; os dois OH formam um \ce{H2O} ($18.015$), \qty{4.32}{\percent}.
(b) \ce{K2O} $47.10/417.25 = \qty{11.3}{\percent}$; \ce{MgO} $3 \times 40.304/417.25 =
\qty{29.0}{\percent}$; \ce{Al2O3} $50.98/417.25 = \qty{12.2}{\percent}$; \ce{SiO2} $3 \times
60.083/417.25 = \qty{43.2}{\percent}$; total \qty{95.7}{\percent}. (c) A microssonda não
detecta o hidrogênio: falta a uma análise de mica sua água, cerca de \qty{4}{\percent}, e seu total deve ficar perto de
96.
\end{solution}

\begin{solution}{exo:b2:field-lab-methods-2:8}
(a) $\beta = \ln(0.1194/0.11905)/\ln(85.909261/87.905612) = 0.002936/(-0.02297) = -0.1278$.
(b) $0.71062 \times (86.908878/85.909261)^{-0.1278} = 0.70957$. (c) A razão 86/88 medida está abaixo da
razão verdadeira: o isótopo pesado foi favorecido, e a 87/86 medida está alta demais.
\end{solution}

\begin{solution}{exo:b2:field-lab-methods-2:9}
(a) $0.1000 \times 15.271 \times 10^{-6}/83.9242 \times 1.3599 \times 10^{-3} = \qty{2.474e-11}{mol}$.
(b) Por \eqref{eq:b2:field-lab-methods-2:id}, $N_x = 2.474 \times 10^{-11} \times (732.8 -
0.874)/(0.874 - 0.0568) = \qty{2.216e-8}{mol}$ de $^{86}\mathrm{Sr}$; estrôncio: $2.216 \times
10^{-8}/0.0986 \times 87.62 = \qty{1.969e-5}{g}$ em \qty{0.0500}{g}, \qty{394}{ppm} (393.9).
(c) Uma recuperação de \qty{99.95}{\percent}: o método é exato dentro do arredondamento da razão.
\end{solution}

\begin{solution}{exo:b2:field-lab-methods-2:10}
(a) $R_m = \sqrt{732.8 \times 0.05680} = 6.45$, onde $\mathcal M = 1.018$. (b) $\mathcal M = 2.31$
em 0.1, 1.13 em 0.5, 1.038 em 1.664 e 1.16 em 100. (c) Perto de $R_x$ a mistura quase não difere
da amostra, perto de $R_s$ do traçador, e a diferença que carrega a informação perde-se no
erro de $R_m$: $\mathcal M$ diverge nas duas extremidades. Entre elas o mínimo é muito achatado
($\mathcal M < 1.2$ de cerca de 0.4 a 100), de modo que o analista só precisa da ordem de grandeza certa
do traçador, que uma concentração aproximada, por FRX por exemplo, fornece.
\end{solution}

\begin{solution}{exo:b2:field-lab-methods-2:11}
(a) $\sigma = \sqrt{0.62 \times 0.38/400} = 0.0243$: $62 \pm 2.4$. (b) $(65 - 62)/2.43 = 1.2$
desvio-padrão: a rocha pode muito bem ser um granodiorito. (c) $\sigma = 0.015$ exige $0.62 \times
0.38/0.015^2 = 1047.1$, ou seja, 1048 pontos de feldspato.
\end{solution}

\begin{solution}{exo:b2:field-lab-methods-2:12}
(a) Cada ponto erra com probabilidade $1 - p$, independentemente: $(1-p)^n$. (b) $0.99^n \le 0.05$
dá $n \ge \ln 0.05/\ln 0.99 = 298.1$: 299 pontos. (c) $n \ge \ln 0.05/\ln 0.9995 = 5990$: uma contagem em malha
é a ferramenta errada; um mineral acessório raro é procurado varrendo a lâmina inteira.
\end{solution}

\begin{solution}{pb:b2:field-lab-methods-2:1}
\noindent\textbf{Parte I.}
\textbf{1.} Quartzo $372/1200 = \qty{31.0}{\percent}$, feldspato alcalino 29.0, plagioclásio 27.0, biotita
9.0, muscovita 3.0, óxidos opacos 1.0.
\textbf{2.} $\sqrt{0.31 \times 0.69/1200} = 0.0134$; muscovita: $\sqrt{0.97/(0.03 \times 1200)} =
\qty{16.4}{\percent}$.
\textbf{3.} $M = 9 + 3 + 1 = \qty{13.0}{\percent}$; $M' = 13 - 3 = \qty{10.0}{\percent}$ sem a
muscovita: leucocrática.
\textbf{4.} $Q + A + P = 1044$ pontos: Q $= 35.6$, A $= 33.3$, P $= 31.0$.
\textbf{5.} $324/672 = \qty{48.2}{\percent}$, com Q entre 20 e 60: um granito, e, entre 35 e
65, um monzogranito.
\textbf{6.} $\sqrt{0.482 \times 0.518/672} = 0.0193$: os limites em 35 e 65 ficam a 6.9 e 8.7
desvios-padrão. O nome é seguro.
\textbf{7.} $0.31 \times 0.69/0.01^2 = 2139$ pontos.

\noindent\textbf{Parte II.}
\textbf{8.} $0.0051/0.8000 = \qty{0.64}{\percent}$.
\textbf{9.} $0.7949/8.7949 = 0.0904$, e $72 \times 0.0904 = \qty{6.51}{\percent}$ de \ce{SiO2} no
vidro: a fusão dilui a amostra onze vezes.
\textbf{10.} $\bar I = 24.91$, $\bar c = 62.283$, $\sum(I - \bar I)^2 = 31.918$, $\sum(I - \bar I)(c -
\bar c) = 80.009$: $a = 2.5067$ por kcps, $b = 62.283 - 2.5067 \times 24.91 = -0.16$.
\textbf{11.} $2.5067 \times 28.63 - 0.16 = \qty{71.61}{\percent}$.
\textbf{12.} BCR-2 $-0.04$, AGV-2 $+0.05$, RGM-1 $-0.01$: todos menores que os desvios-padrão
do método (0.1 a 0.2), de modo que a reta é adequada nessa faixa.
\textbf{13.} $71.61 + 27.15 + 0.64 = \qty{99.40}{\percent}$, entre 97 e 101: não é repetida.
\textbf{14.} $71.61 \times 100/(71.61 + 27.15) = \qty{72.51}{\percent}$.

\noindent\textbf{Parte III.}
\textbf{15.} $0.0986/0.8258 = 0.1194$: o valor de normalização é a razão natural.
\textbf{16.} $\beta = 0.007313/(-0.022972) = -0.3183$.
\textbf{17.} $0.71655 \times (86.908878/85.909261)^{-0.3183} = 0.71392$.
\textbf{18.} A razão 86/88 medida está a $0.11853/0.1194 - 1 = \qty{-0.73}{\percent}$ do valor verdadeiro
em duas unidades de massa: cerca de \qty{0.36}{\percent} por unidade de massa, a favor dos isótopos pesados.
\textbf{19.} A razão $^{86}\mathrm{Sr}/^{88}\mathrm{Sr}$ do traçador é 0.7661: a razão 86/88 da mistura já não é
a natural, 0.1194, e fixá-la em 0.1194 trataria o traçador como um viés instrumental.

\noindent\textbf{Parte IV.}
\textbf{20.} Soma das razões a 88: $561.4 + 0.7661 + 0.1819 + 1 = 563.348$; $561.4/563.348 =
\qty{99.65}{\percent}$ e $0.7661/563.348 = \qty{0.13599}{\percent}$.
\textbf{21.} $0.2500 \times 15.271 \times 10^{-6}/83.9242 = \qty{4.549e-8}{mol}$ de estrôncio, e
$\times 1.3599 \times 10^{-3} = \qty{6.186e-11}{mol}$ de $^{86}\mathrm{Sr}$.
\textbf{22.} $561.4/0.7661 = 732.8$ e $0.0056/0.0986 = 0.05680$.
\textbf{23.} $N_x = 6.186 \times 10^{-11} \times (732.8 - 1.6812)/(1.6812 - 0.0568) =
\qty{2.784e-8}{mol}$.
\textbf{24.} $\mathcal M = 1.6812 \times 732.74/(731.12 \times 1.6244) = 1.04$; a razão acrescenta
$1.04 \times \qty{0.01}{\percent}$, o traçador $0.031/15.271 = \qty{0.20}{\percent}$; em quadratura,
\qty{0.20}{\percent}, dominados pela calibração do traçador.
\textbf{25.} Estrôncio: $2.784 \times 10^{-8}/0.0986 = \qty{2.824e-7}{mol}$, isto é, $2.824 \times
10^{-7} \times 87.62 = \qty{2.474e-5}{g}$ em \qty{0.1000}{g}. O granito de Varden contém
\textbf{\qty[text-series-to-math=true]{247}{ppm}} de estrôncio, $\pm 0.5$.
\end{solution}
